% 量子力学1 量子力学1-10 問題6（031-expand。当時は 量子力学1-4 問題6・kai04-l2-expand） 解いた日 2026-10-28
%% kotae
\toi{1}{◎}{}
\begin{align*}
\int_{0}^{L}\abs{f}^{2}\dd x&=\frac{30}{L^{5}}\int_{0}^{L}x^{2}(L-x)^{2}\dd x\\
&=\frac{30}{L^{5}}\int_{0}^{L}(x^{2}L^{2}-2x^{3}L+x^{4})\dd x\\
&=\frac{30}{L^{5}}\Bigl[\frac{1}{3}x^{3}L^{2}-\frac{1}{2}x^{4}L+\frac{1}{5}x^{5}\Bigr]_{0}^{L}\\
&=\frac{30}{L^{5}}\cdot L^{5}\Bigl(\frac{1}{3}-\frac{1}{2}+\frac{1}{5}\Bigr)=1
\end{align*}
\toi{2}{◎}{}
\[ f(x)\chi_{m}^{*}(x)=\sum_{n}c_{n}\chi_{n}(x)\chi_{m}^{*}(x) \]
\begin{align*}
\int_{0}^{L}f\chi_{m}^{*}\dd x&=\int_{0}^{L}\sum_{n}c_{n}\chi_{m}^{*}\chi_{n}\dd x\\
&=\sum_{n}\int_{0}^{L}c_{n}\chi_{m}^{*}\chi_{n}\dd x\\
&=\sum_{n}c_{n}\delta_{nm}\\
&=c_{m}
\end{align*}
\[ \therefore\ c_{n}=\int_{0}^{L}\chi_{n}^{*}(x)f(x)\dd x \]
\toi{3}{◎}{}
\begin{align*}
c_{n}&=\int_{0}^{L}\sqrt{\frac{2}{L}}\sin\frac{n\pi x}{L}\sqrt{\frac{30}{L^{5}}}\,x(L-x)\dd x\\
&=\frac{2\sqrt{15}}{L^{3}}\int_{0}^{L}(-x^{2}+Lx)\sin\frac{n\pi x}{L}\dd x
\end{align*}
\begin{align*}
&\int_{0}^{L}x\sin\frac{n\pi x}{L}\dd x\\
&=\Bigl[\frac{-Lx}{n\pi}\cos\frac{n\pi x}{L}\Bigr]_{0}^{L}+\int_{0}^{L}\frac{L}{n\pi}\cos\frac{n\pi x}{L}\dd x\\
&=\frac{-L^{2}}{n\pi}\cos n\pi+\Bigl(\frac{L}{n\pi}\Bigr)^{2}\Bigl[\sin\frac{n\pi x}{L}\Bigr]_{0}^{L}\\
&=\frac{-L^{2}}{n\pi}(-1)^{n}
\end{align*}
\begin{align*}
&\int_{0}^{L}x^{2}\sin\frac{n\pi x}{L}\dd x\\
&=\Bigl[\frac{-Lx^{2}}{n\pi}\cos\frac{n\pi x}{L}\Bigr]_{0}^{L}+\int_{0}^{L}\frac{2Lx}{n\pi}\cos\frac{n\pi x}{L}\dd x\\
&=\frac{-L^{3}}{n\pi}(-1)^{n}+\frac{2L}{n\pi}\int_{0}^{L}x\cos\frac{n\pi x}{L}\dd x
\end{align*}
\begin{align*}
&\int_{0}^{L}x\cos\frac{n\pi x}{L}\dd x\\
&=\Bigl[\frac{Lx}{n\pi}\sin\frac{n\pi x}{L}\Bigr]_{0}^{L}-\int_{0}^{L}\frac{L}{n\pi}\sin\frac{n\pi x}{L}\dd x\\
&=+\Bigl(\frac{L}{n\pi}\Bigr)^{2}\Bigl[\cos\frac{n\pi x}{L}\Bigr]_{0}^{L}\\
&=\Bigl(\frac{L}{n\pi}\Bigr)^{2}\bigl((-1)^{n}-1\bigr)
\end{align*}
\[ =\frac{-L^{3}}{n\pi}(-1)^{n}+2\Bigl(\frac{L}{n\pi}\Bigr)^{3}\bigl((-1)^{n}-1\bigr) \]
\begin{align*}
c_{n}&=\frac{2\sqrt{15}}{L^{3}}\biggl(\frac{L^{3}}{n\pi}(-1)^{n}-2\Bigl(\frac{L}{n\pi}\Bigr)^{3}\bigl((-1)^{n}-1\bigr)\\
&\qquad-\frac{L^{3}}{n\pi}(-1)^{n}\biggr)\\
&=4\sqrt{15}\,\frac{1}{(n\pi)^{3}}\bigl(1-(-1)^{n}\bigr)
\end{align*}
奇数：$\dfrac{8\sqrt{15}}{(n\pi)^{3}}$，偶数：$0$
\toi{4}{◎}{}
\[ \abs{c_{1}}^{2}=\frac{64\cdot15}{\pi^{6}}=0.9986 \]
\[ \abs{f}^{2}=\Bigl(\sum_{n}c_{n}\chi_{n}\Bigr)^{*}\Bigl(\sum_{m}c_{m}\chi_{m}\Bigr) \]
\begin{align*}
\int\abs{f}^{2}\dd x&=\sum_{n}\sum_{m}c_{n}^{*}c_{m}\int\chi_{n}^{*}\kakimachigai{\chi_{n}}{\chi_{m}}\dd x\\
&=\sum_{n}\sum_{m}c_{n}^{*}c_{m}\delta_{nm}\\
&=\sum_{n}\abs{c_{n}}^{2}
\end{align*}
\begin{align*}
1&=\sum_{n\,\text{奇}}\frac{960}{(n\pi)^{6}}\\
&=\frac{960}{\pi^{6}}\sum_{n\,\text{奇}}\frac{1}{n^{6}}=\frac{960}{\pi^{6}}\cdot\frac{\pi^{6}}{960}=1
\end{align*}
%% kizuki
